\chapter{Propiedades Mecánicas y Comportamiento Viscoelástico}
\label{cap:mecanicas}

Los polímeros poseen un comportamiento mecánico único denominado viscoelasticidad. A diferencia de los sólidos puramente elásticos (que obedecen la ley de Hooke) y de los líquidos puramente viscosos (que obedecen la ley de Newton para fluidos), los polímeros exhiben ambas propiedades simultáneamente debido a la flexibilidad y enredo físico de sus largas cadenas bajo esfuerzos deformantes mecánicos.

\section{Modelos Viscoelásticos Clásicos}

Para aproximar cuantitativamente el comportamiento viscoelástico lineal, se recurre a modelos analógicos mecánicos compuestos por combinaciones de resortes elásticos ideales de constante $E$ y amortiguadores viscosos ideales de viscosidad $\eta$:

\subsection{Modelo de Maxwell (Comportamiento en Relajación)}
Consiste en un resorte y un amortiguador acoplados en serie. Su ecuación constitutiva diferencial general se define como:
\begin{equation}
    \frac{d\epsilon}{dt} = \frac{1}{E} \frac{d\sigma}{dt} + \frac{\sigma}{\eta}
    \label{eq:maxwell}
\end{equation}
Si sometemos el sistema a una deformación constante e instantánea ($\epsilon = \epsilon_0$ y $\frac{d\epsilon}{dt} = 0$), el esfuerzo $\sigma(t)$ se relaja de manera exponencial a través del tiempo siguiendo la siguiente expresión matemática:
\begin{equation}
    \sigma(t) = E \epsilon_0 e^{-t/\tau} \quad \text{donde} \quad \tau = \frac{\eta}{E}
    \label{eq:relajacion_maxwell}
\end{equation}
donde $\tau$ representa el tiempo de relajación característico del material polimérico.

\subsection{Modelo de Kelvin-Voigt (Comportamiento en Fluencia/Creep)}
Compuesto por un resorte y un amortiguador colocados en paralelo. Su ecuación constitutiva se define mediante la siguiente relación:
\begin{equation}
    \sigma = E\epsilon + \eta \frac{d\epsilon}{dt}
    \label{eq:kelvin_voigt}
\end{equation}
Si se aplica un esfuerzo constante e instantáneo $\sigma_0$ (ensayo de fluencia), el material sufre una deformación progresiva retardada que se estabiliza asintóticamente hacia el módulo elástico del resorte.

\section{Análisis Mecánico Dinámico (DMA)}

El comportamiento bajo esfuerzos oscilatorios dinámicos alternantes revela la respuesta elástica y viscosa en función de la temperatura y de la frecuencia. En el DMA se definen las siguientes variables:
\begin{description}
    \item[Módulo de Almacenamiento ($E'$):] Medida de la energía mecánica almacenada elásticamente de forma reversible por el polímero.
    \item[Módulo de Pérdida ($E''$):] Medida de la energía disipada irreversiblemente en forma de calor debido al flujo viscoso interno de las cadenas.
    \item[Factor de Pérdida ($\tan \delta$):] Razón de amortiguamiento, indicando la amortiguación del material frente a vibraciones mecánicas:
        \begin{equation}
            \tan \delta = \frac{E''}{E'}
            \label{eq:tan_delta}
        \end{equation}
\end{description}
Un pico muy pronunciado en la curva de $\tan \delta$ en función de la temperatura define experimentalmente la temperatura de transición vítrea ($T_g$) de la muestra.
